% ------------------------------------------------------------------------
%
\begin{frame}{A mini functional language}

  Rewrite systems are a good mathematical and mental model for
  understanding functional languages.

  \bigskip

  It is time now to learn an actual functional language, so we can run
  programs. We are going to start with a mini language that is
  actually a subset of \OCaml, then grow it step by step.

  \bigskip

  Within rewrite systems, a `program' is simply a term that is reduced
  until reaching a stuck term. While this is a good theoretic
  framework to think about computation in general, it is not conducive
  to programming in the large, nor understanding the behaviour of
  programs.

  \bigskip

  We want programs to be made of different parts that can be reused in
  different contexts, for example, variables should be defined
  somewhere and reused at some another points.

\end{frame}

% ------------------------------------------------------------------------
%
\begin{frame}{Growing a language}

  This presentation is principled, that is, it relies on a systematic
  classification of the \textbf{syntactic constructs} and their
  \textbf{formal semantics}.

  \bigskip

  When we reach a certain complexity, we will drop the formal
  descriptions.

  \bigskip

  Both theoretical and pragmatic, our approach is complementary to
  those usually found in books for general audiences.

  \bigskip

  Let mini-ML be the name of the subset of \OCaml we are going to
  grow.

\end{frame}

% ------------------------------------------------------------------------
%
\begin{frame}{Appending a stack}

  The notation we introduced for stacks in the context of rewrite
  systems actually follows the \OCaml convention.

  \bigskip

  In the context of functional programming, stacks are often called
  \textbf{lists}.

  \bigskip

  Here are different ways to express the same list in \OCaml:
  \begin{equation*}
    1::(2::3::(4::\el)) = 1::2::3::4::\el = [1;2;3;4].
  \end{equation*}

  Before starting, let us see how a familiar function defined with a
  rewrite system looks like in \OCaml.

  \bigskip

  Let us recall the appending of a stack:
  \begin{align*}
    \fun{cat}(        \el,t) &\smashedrightarrow{\alpha} t;\\
    \fun{cat}(\cons{x}{s},t) &\smashedrightarrow{\beta} \cons{x}{\fun{cat}(s,t)}.
  \end{align*}

\end{frame}

% ------------------------------------------------------------------------
%
\begin{frame}[fragile]{Appending a stack}

  In \OCaml, the function name is not repeated (once for each rule).

  \bigskip

  Instead, it is introduced once by a \textbf{binder} called `\Xlet',
  like so:
  \vspace*{-3mm}
\begin{alltt}
\small
\textbf{let} cat = \textbf{function} ...
\end{alltt}
  The complete definition in \OCaml is:
  \vspace*{-3mm}
\begin{alltt}
\small
\textbf{let} \textbf{rec} cat = \textbf{function}
\ \ \ [], t \(\rightarrow\) t
\(\mid\) \ttcons{x}{s}, t \(\rightarrow\) \ttcons{x}{\fun{cat}(s,t)}
\end{alltt}

   Note the addition of the keyword \Xrec: this is meant so the call
   \(\fun{cat}(s,t)\) indeed refers to the function being defined.

   \bigskip

   With rewrite systems, recursion is not an issue: any call matching
   a pattern can be rewritten.

   \bigskip

   Note also how the two rules are introduced by the keyword
   \Xfunction and are separated by a vertical line.

\end{frame}

% ------------------------------------------------------------------------
%
\begin{frame}{Sentences}

  An \OCaml program is a sequence of sentences. A \textbf{sentence},
  or \textbf{global definition}, is defined by the following cases,
  where~\(e\) denotes an expression, \(x\)~is a variable, and `\Xlet'
  and~`\Xrec' are keywords:

  \medskip

  \begin{tabular}{rll}
    $\bullet$ & \emph{global definition}
    & \phrase{$\Xlet \; x \; = \; e$}\\
    $\bullet$ & \emph{global recursive definition}
    & \phrase{$\Xlet \; \Xrec \; x \; = \; e$}
  \end{tabular}

  \medskip

  Global definitions can be optionally followed by two semicolons:
  these are necessary when inputting the sentences in the toplevel
  loop (prompted after running the command \texttt{\small ocaml} from
  a Unix shell).

  \bigskip

  Note that the keyword `\Xrec' is necessary to enable recursion. (We
  will see later how that works precisely.)

  \bigskip

  For~\(e\) to be an expression means that \(e\)~denotes a part of a
  sentence which is classified as an expression according to its
  \textbf{syntax}.

\end{frame}

% ------------------------------------------------------------------------
%
\begin{frame}[fragile]{Metavariables}

  \begin{itemize}

    \item We say that \(e\)~is a \textbf{metavariable} because, being
      a name, it is a variable, but that variable does not belong to
      the language being described (the \textbf{object language},
      here, \OCaml) and, instead, exists only in the descriptive
      language (the \textbf{metalanguage}).

    \item In the following
      \begin{center}
        \phrase{$\Xlet \; x \; = \; e$}
      \end{center}
      the metavariable~\(x\) (mind the italics) denotes an infinite
      set of \OCaml variables and we must not confuse it with the
      \OCaml variable~\ident{x} (no italics). Similarly, the
      metavariable \(e\)~denotes an infinity of expressions, as in the
      following:
      \vspace*{-3mm}
\begin{alltt}
\small
\textbf{let} e = 3
\textbf{let} x = e
\textbf{let} a = x + e
\end{alltt}

  \end{itemize}

\end{frame}

% ------------------------------------------------------------------------
%
\begin{frame}{Expressions}

  \begin{itemize}

    \item An expression \(e\)~is \textbf{recursively defined by} the
      following \textbf{cases}:

      \smallskip

      \begin{tabular}{@{\,}r@{\,\,}ll}
        $\bullet$
        & \emph{variable}
        & my\_var, x, y, z, etc. \\
        $\bullet$
        & \emph{unit or integer constant}
        & \unit \ or \textsf{0} or \textsf{1} or \textsf{2}, etc.\\
        $\bullet$
        & \emph{function} (or \emph{abstraction})
        & $\Xfun \; x \rightarrow e$\\
        $\bullet$
        & \emph{call} (or \emph{application})
        & $e_1 \; e_2$ \\
        $\bullet$
        & \emph{arithmetic operator}
        & \lpar\texttt{+}\rpar \ \lpar\texttt{-}\rpar \ \lpar\texttt{/}\rpar
        \ \lpar\texttt{*}\rpar\\
        $\bullet$
        & \emph{arithmetic operation}
        & $e_1$ \texttt{+} $e_2$ or $e_1$ \texttt{-} $e_2$,
        etc.\\
        $\bullet$
        & \emph{parenthesis}
        & $\lpar{e}\rpar$\\
        $\bullet$
        & \emph{local definition}
        & $\Xlet \; x = e_1 \; \Xin \; e_2$
      \end{tabular}

    \item Note how the keywords `\Xrec' and `\Xfunction', as well as
      the definition of multiple cases (as seen in the definition of
      `\fun{cat}' earlier) is not found above: this is because we want
      to focus on other aspects of \OCaml first, like
      higher\hyp{}order functions.

    \item Instead, we have `\Xfun', which enables only one case made
      of a variable (more generally, irrefutable patterns).

  \end{itemize}

\end{frame}

% ------------------------------------------------------------------------
%
\begin{frame}[fragile]{Examples of expressions}

  \begin{itemize}

    \item Note that what we print nicely as `\(\rightarrow\)' is
      actually written `\texttt{->}' in the source code.

    \item Here are examples of \OCaml sentences making up a program,
      step by step.

    \item First, a numerical constant:
      \vspace*{-3mm}
\begin{alltt}
\small
\textbf{let} x = 0
\end{alltt}

    \item Note that constants can be defined by means of a rewrite
      system as functions taking the empty tuple:
      \begin{equation*}
        \fun{x}() \rightarrow 0
      \end{equation*}

    \item The next sentence is
      \vspace*{-3mm}
\begin{alltt}
\small
\textbf{let} id = \textbf{fun} x \(\rightarrow\) x
\end{alltt}

    \item This is equivalent to the rewrite system
      \begin{equation*}
        \fun{id}(x) \rightarrow x
      \end{equation*}

  \end{itemize}

\end{frame}

% ------------------------------------------------------------------------
%
\begin{frame}[fragile]{Examples of expressions}

  \begin{itemize}

    \item Next we have another constant being defined by binding a
      variable to the value of a call:
      \vspace*{-3mm}
\begin{alltt}
\small
\textbf{let} y = 2 \textbf{in} id y
\end{alltt}
      This is equivalent to the call \texttt{\small id(2)}.

    \item Now
      \vspace*{-3mm}
\begin{alltt}
\small
\textbf{let} x = (\textbf{fun} x \(\rightarrow\) \textbf{fun} y \(\rightarrow\) x + y) 1 2
\end{alltt}
      is equivalent to
      \vspace*{-3mm}
\begin{alltt}
\small
\textbf{let} x = 1 + 2
\end{alltt}
      (We will see why later.)

    \item Finally
      \vspace*{-3mm}
\begin{alltt}
\small
\textbf{let} z = x + 1
\end{alltt}
       is obvious.

  \end{itemize}

\end{frame}

% ------------------------------------------------------------------------
%
\begin{frame}{Functions}

  \begin{itemize}

    \item Variables must start with a lowercase character.

    \item The arrow is right-associative, so the expression
      \begin{equation*}
        \Xfun \; x_1 \rightarrow \Xfun \; x_2
        \rightarrow \ldots
        \rightarrow \Xfun \; x_n \rightarrow e
      \end{equation*}
      is equivalent to \(\Xfun \; x_1 \rightarrow (\Xfun
      \; x_2 \rightarrow (\ldots \rightarrow
      (\Xfun \; x_n \rightarrow e)) \; \ldots)\).

    \item Function calls are left\hyp{}associative, so
      \begin{equation*}
        e_1 \; e_2 \; e_3 \; \ldots \; e_n \;=\; (((\ldots
        (e_1 \; e_2) \; e_3) \ldots) \; e_n).
      \end{equation*}

    \item Function calls have higher priority than operator calls, for
      example,
      \begin{equation*}
        f \; 3 + 4 \; = \; (f \; 3) + 4 \; = \; f(3) + 4.
      \end{equation*}

    \item Operator calls have higher syntactic priority than
      abstractions, for instance,
      \begin{equation*}
        \Xfun \; x \rightarrow x + y \quad\text{is the same as}\quad
        \Xfun \; x \rightarrow (x + y).
      \end{equation*}

  \end{itemize}

\end{frame}

% ------------------------------------------------------------------------
%
\begin{frame}[fragile]{Function calls}

  Note the difference between \texttt{\small f x y} and \texttt{\small
    f(x,y)}.

  \bigskip

  The former is to be understood as \texttt{\small (f(x))(y)} which
  means that there are \emph{two} calls:
  \begin{enumerate}

  \item First, the call \texttt{\small f(x)}
    is evaluated. Let us assume that its value is \texttt{\small g}.

  \item Second, the call \texttt{\small g(y)} is evaluated.

  \end{enumerate}

  In other words, the first case is equivalent to
  \vspace*{-3mm}
\begin{alltt}
\small
\textbf{let} g = f x \textbf{in} g y
\end{alltt}

\end{frame}

% ------------------------------------------------------------------------
%
\begin{frame}{Abstract syntax}

  A graphical representation of programs (here, expressions) as trees
  proves very handy when studying their semantics (meaning):

  \begin{center}
    \begin{tabular}{lc}
      \toprule
      \multicolumn{1}{c}{Expression}
        & \multicolumn{1}{c}{Tree}\\
        \toprule
        $x$
        & $x$\\
        $\Xfun \; x \rightarrow e$
        & \includegraphics[bb=148 638 171 662]{fun_tree}\\
        $e_1 \; e_2$
        & \includegraphics[bb=148 637 177 663]{app_tree}\\
        \bottomrule
    \end{tabular}
  \end{center}

  Note that we needed some symbol for the root of the application tree
  (\texttt{\small \$}). Traditionally, that symbol is (\texttt{\small
    @}), but \OCaml uses it already for something else.

  \bigskip

  When trees represent programs, they are called \textbf{abstract
    syntax trees} (AST).

\end{frame}

% ------------------------------------------------------------------------
%
\begin{frame}{Abstract syntax}

  Intuitively, the method for constructing abstract syntax trees
  consists in fully parenthesising the expression. Each parenthesis
  corresponds to a subexpression and each subexpression corresponds to
  a subtree.

  \bigskip

  The tree is built from the root, down to the leaves, by traversing
  subexpressions from the outermost to the innermost, that is, the
  most embedded ones. For instance, consider
  \begin{center}
    \includegraphics[bb=148 617 233 660]{arith_tree1}
    \qquad
    \includegraphics[bb=148 599 246 662]{arith_tree2}
  \end{center}
  \centerline{\texttt{\small ((1+2)*(5/1)) \quad\qquad \texttt{\small \textbf{let} x = 1 \textbf{in} (1+2)*5/1}}}

\end{frame}

% ------------------------------------------------------------------------
%
\begin{frame}[fragile]{Abstract syntax}

  Another example:
    \vspace*{-3mm}
\begin{alltt}
\small
\textbf{let} x = 1 \textbf{in} ((\textbf{let} x = 2 \textbf{in} x) + x)
\end{alltt}
  corresponds to the abstract syntax tree
  \begin{center}
    \includegraphics[bb=148 599 221 662]{arith_tree3}
  \end{center}

  This example features two local definitions of `x': is one the
  redefinition of the other? Is one hiding the other?

  \bigskip

  There are also two uses of `x': which use corresponds to which
  definition?

\end{frame}

% ------------------------------------------------------------------------
%
\begin{frame}{Bindings, environments, scope}

  To answer theses questions, we must first understand
  \begin{itemize}

    \item what a \textbf{binding} is,

    \item what the \textbf{scope} of a binding is, and

    \item what an \textbf{environment} is.
  \end{itemize}

  The expression we called \textbf{local definition} earlier:
  \begin{equation*}
    \Xlet \; x = e_1 \; \Xin \; e_2
  \end{equation*}
  is evaluated as follows:
  \begin{enumerate}

    \item The expression \(e_1\)~is evaluated,

    \item assuming it has a value~\(v_1\), then \(v_1\) is
      \textbf{bound} to the variable~\(x\),

    \item the expression \(e_2\)~is evaluated under the assumption
      that \(x\)~is bound to~\(v_1\).

  \end{enumerate}

\end{frame}

% ------------------------------------------------------------------------
%
\begin{frame}{Bindings, environments, scope}

  A local definition
  \begin{equation*}
    \Xlet \; x = e_1 \; \Xin \; e_2
  \end{equation*}
  associates the value~\(v_1\) of expression~\(e_1\) to the variable
  denoted by~\(x\), making up a \textbf{binding} noted \(x \; \mapsto
  \; v_1\), and then proceeds to evaluate~\(e_2\) with this binding
  available.

  \bigskip

  The binding \(x \mapsto v_1\) is not usable outside the evaluation
  of~\(e_2\). We say that the \textbf{scope} of the definition is made
  of the subtrees of the abstract syntax tree corresponding to~\(e_2\)
  where the binding of that definition is usable (`visible').

  \bigskip

  A mapping from variables to values is called an
  \textbf{environment}. Let us note \(\rho_\varnothing\) the empty
  environment. We write
  \begin{equation*}
    (x \mapsto v) \oplus \rho \quad\text{or, simply,}\quad x \mapsto v
    \oplus \rho,
  \end{equation*}
  to mean that the binding \(x \mapsto v\) was added to the
  environment~\(\rho\).

\end{frame}

% ------------------------------------------------------------------------
%
\begin{frame}[fragile]{Bindings, environments, scope}

  Environments are functions and are updated as follows:
  \begin{equation}
    (x \mapsto v \oplus \rho )(y) \triangleq
    \begin{cases}
      v         & \text{if} \; x = y,\\
      \sigma(y) & \text{otherwise.}
    \end{cases}
  \end{equation}
  with \(\forall x.\rho_\varnothing(x) = \bot\) and \(\bot\) is
  interpreted as an \textbf{error}. Errors are special
  values. Consider again the program
  \vspace*{-3mm}
\begin{alltt}
\small
\textbf{let} x = 0 \textbf{in}
\textbf{let} id = \textbf{fun} x \(\rightarrow\) x \textbf{in}
\textbf{let} y = id x \textbf{in}
\textbf{let} x = (\textbf{fun} x \(\rightarrow\) \textbf{fun} y \(\rightarrow\) x + y) 1 2
\textbf{in} x + 1
\end{alltt}

  The variable~\texttt{\small x} in the expression \texttt{\small x+1}
  denotes the value of the expression bound to the
  variable~\texttt{\small x} in the previous line, not the first.

  \bigskip

  Bindings are ordered in the environment \textbf{by the order of
    their definitions}. Let us follow the evolution of the environment.

\end{frame}

% ------------------------------------------------------------------------
%
\begin{frame}{Bindings, environments, scope}

  \begin{enumerate}

    \item The environment is initially empty: \(\rho_\varnothing\);

    \item \texttt{\small \textbf{let} x = 0 \textbf{in}
      \ldots}\\ yields $(\ident{x} \mapsto 0) \oplus
      \rho_\varnothing$;

    \item \texttt{\small \textbf{let} id = \textbf{fun} x
      $\rightarrow$ x \textbf{in} \ldots}\\ yields $(\ident{id}
      \mapsto \Xfun \; \ident{x} \rightarrow \ident{x}) \oplus
      (\ident{x} \mapsto 0) \oplus \rho_\varnothing$;

    \item \texttt{\small \textbf{let} y = id x \textbf{in}
      \ldots}\\ yields $(\ident{y} \mapsto 0) \oplus (\ident{id}
      \mapsto \Xfun \; \ident{x} \rightarrow \ident{x}) \oplus
      (\ident{x} \mapsto 0) \oplus \rho_\varnothing$;

    \item \texttt{\small \textbf{let} x = \ldots}\\ yields
      $(\underline{\ident{x} \mapsto 3}) \oplus (\ident{y}
      \mapsto 0) \oplus (\ident{id} \mapsto \Xfun \;
      \ident{x} \rightarrow \ident{x}) \oplus (\ident{x} \mapsto
      0) \oplus \rho_\varnothing$.

    \item The expression \textsf{\small x+1} is evaluated in \(3+1=4\)
      because \(\textsf{x} \mapsto 3\) (the last binding)
      \textbf{hides} \(\textsf{x} \mapsto 0\) (the first
      binding).

    \item In other words, the first definition of~\texttt{\small x} was
      \textbf{out of scope} in \texttt{\small x+1}.

  \end{enumerate}

\end{frame}

% ------------------------------------------------------------------------
%
\begin{frame}{Free variables}

  \begin{itemize}

    \item If a variable~\(x\) is not bound to a value in an
      environment~\(\rho\), we say that \(x\)~is \textbf{free}
      in~\(\rho\).

    \item By definition, a variable is \textbf{free} in an expression
      if it is not bound by a `\Xlet' or a `\Xfun'.

    \item We can understand this concept on a graphic representation
      of the abstract syntax tree.

    \item From each variable occurrence, we move up, towards the
      root.

    \item If we encounter a `\Xlet' that binds that variable (in the
      left child), we create an oriented edge from the variable to
      that `\Xlet', and we stop: the variable is bound.

    \item Otherwise, if we reach the root (no such `\Xlet' has been
      found), then the variable is free.

  \end{itemize}

\end{frame}

% ------------------------------------------------------------------------
%
\begin{frame}{Free variables}

  \begin{figure}
    \begin{center}
      \includegraphics[bb=148 599 221 662]{free_var1}
      \caption{\textbf{let} x = 1 \textbf{in}
      ((\textbf{let} x = 2 \textbf{in} x) + x)}
    \end{center}
  \end{figure}

  \begin{figure}
    \begin{center}
      \includegraphics[bb=148 582 211 662]{free_var2}
    \end{center}
    \caption{\textbf{fun} y \(\rightarrow\) x + (\textbf{fun} x
      \(\rightarrow\) x) y}
  \end{figure}

\end{frame}

% ------------------------------------------------------------------------
%
\begin{frame}{Free variables}

  A similar situation arises with functions: `\Xfun' is a binder, just
  like `\Xlet'.

  \bigskip

  In \Xfun $x \rightarrow e$, the variable~\(x\) (called a
  \textbf{parameter}) may shadow (hide) another variable~\(x\) defined
  on the path to the root.

  \bigskip

  See previous page for two examples.

  \bigskip

  We note \(\FV{e}\) the set of the variables free in~\(e\).

\end{frame}

% ------------------------------------------------------------------------
%
\begin{frame}{Free variables}

  It is easy to formally define \(\mathcal{F}\) by case and
  recursively:
  \begin{align*}
    \FV{x} &= \{x\}, \quad\text{where $x$ is a  variable}\\
    \FV{c} &= \varnothing, \quad\text{where $c$ is a constant}\\
    \FV{\Xfun \; x \rightarrow e} &=
    \FV{e} \backslash \{x\}\\
    \FV{e_1 \; e_2} &= \FV{e_1} \;\cup\; \FV{e_2}\\
    \FV{\Xlet \; x  = e_1 \; \Xin \; e_2} &=
    \FV{e_1} \;\cup\; \FV{e_2} \backslash \{x\}\\
    \FV{e_1 \; \texttt{+} \; e_2} &= \FV{e_1} \;\cup\; \FV{e_2}
  \end{align*}

  Now, instead of drawing abstract syntax trees, let us use these
  equations to compute the free variables of our previous examples.

\end{frame}

% ------------------------------------------------------------------------
%
\begin{frame}{Free variables}

  \begin{equation*}
    \begin{array}{lcl}
      \multicolumn{3}{l}{
        \FV{\Xlet \; \textsf{x} = 1 \;
          \Xin \; (\Xlet \; \textsf{x} = 2 \;
          \Xin \; \textsf{x}) \; + \; \textsf{x}}} \\
      \phantom{XXX} &=& \FV{1} \cup \FV{(\Xlet \; \textsf{x}
        = 2 \; \Xin \; \textsf{x}) +
        \textsf{x}} \backslash \{\textsf{x}\}\\
      \phantom{XXX} &=& \varnothing \cup (\FV{\Xlet \;
        \textsf{x} = 2 \; \Xin \;
        \textsf{x}} \cup \FV{\textsf{x}}) \backslash \{\textsf{x}\}\\
      \phantom{XXX} &=& (\FV{2} \cup \FV{\textsf{x}} \backslash
      \{\textsf{x}\}  \cup \FV{\textsf{x}}) \backslash
      \{\textsf{x}\}\\
      \phantom{XXX} &=& \FV{\textsf{x}} \backslash \{\textsf{x}\}\\
      \phantom{XXX} &=& \varnothing.
    \end{array}
  \end{equation*}

  \bigskip

  \begin{equation*}
    \begin{array}{lcl}
      \multicolumn{3}{l}{ \FV{\Xfun \; \textsf{y} \rightarrow
          \textsf{x} + (\Xfun \; \textsf{x}
          \rightarrow \textsf{x}) \; \textsf{y}}}\\ \phantom{X}
      &=& \FV{\textsf{x} + (\Xfun \; \textsf{x} \rightarrow
        \textsf{x}) \; \textsf{y}} \backslash
      \{\textsf{y}\}\\ \phantom{X} &=& (\FV{\textsf{x}} \cup
      \FV{(\Xfun \; \textsf{x} \rightarrow \textsf{x})
        \; \textsf{y}}) \backslash \{\textsf{y}\}\\ \phantom{X} &=&
      (\{\textsf{x}\} \cup \FV{\Xfun \; \textsf{x}
        \rightarrow \textsf{x}} \cup \FV{\textsf{y}})
      \backslash \{\textsf{y}\}\\ \phantom{X} &=& (\{\textsf{x}\}
      \cup \FV{\textsf{x}} \backslash \{\textsf{x}\} \cup
      \{\textsf{y}\}) \backslash \{\textsf{y}\}\\ \phantom{XXX} &=&
      \{\textsf{x}, \textsf{y}\} \backslash
      \{\textsf{y}\}\\ \phantom{X} &=& \{\textsf{x}\}.
    \end{array}
  \end{equation*}

\end{frame}

% ------------------------------------------------------------------------
%
\begin{frame}{Closed expressions}

  \begin{itemize}

    \item A \textbf{closed} expression is an expression without free
      variables.

    \item That is, an expression~\(e\) is closed if, and only if,
      \(\FV{e} = \varnothing\).

    \item By extension, we say that~\(e\) is \textbf{closed in an
      environment}~\(\rho\) if \(\FV{e} \subseteq \dom{\rho}\) (the
      free variables are bound in~\(\rho\)).

    \item Only a closed expression can be evaluated (executed). That
      is why the first static analysis performed by compilers consists
      in determining the variables which are free in expressions: if
      the program is not closed, it is rejected.

    \item In the case of
      \begin{center}
        \Xfun y $\rightarrow$ x + (\Xfun x $\rightarrow$ x) y,
      \end{center}
      the \OCaml compiler prints out
      \begin{center}
        \texttt{\small \textbf{Unbound value x}}
      \end{center}
      and stops. It is useful that this open expression be rejected at
      compile\hyp{}time and does not cause a run\hyp{}time error.

  \end{itemize}

\end{frame}

% ------------------------------------------------------------------------
%
\begin{frame}{Evaluation}

  \textbf{Evaluation} is a partial function from expressions to
  values. An \textbf{interpreter} is a program implementing an
  evaluation. The partiality models the fact that an evaluation may
  not terminate or end in an error.

  \bigskip

  With rewrite systems, values are defined as stuck expressions, that
  is, expressions that cannot be rewritten further.

  \bigskip

  Here, since we want to define expressions (terms) whose value we can
  reuse somewhere else (in another term), we need the notion of
  \textbf{environment}.

  \bigskip

  That implies that a formal account of evaluation (the equivalent for
  programs of rewriting terms) needs to take into account
  environments.

\end{frame}

% ------------------------------------------------------------------------
%
\begin{frame}{Semantic values}

  This goes even further because we want Mini-ML to be
  \textbf{higher\hyp{}order}, that is, we want functions to be values
  and this begets the question of the meaning of the free variables in
  the function bodies: after passing around a function as argument to
  other functions, what are the values of those variables when the
  threaded function is finally applied?

  \bigskip

  The most common answer is \textbf{lexical (or static) scoping},
  whereby the free variables in a body remain bound to their values at
  the point the function was defined.

  \bigskip

  This means that values corresponding to functions have to carry
  around the environment at the point of their definition.

  \bigskip

  One common way to achieve this is to have a notion of
  \textbf{semantic value}, which differs from \textbf{syntactic
    value}. For example, the mathematical integers are part of the
  semantic values of Mini\hyp{}ML, but not of syntactical values.

\end{frame}

% ------------------------------------------------------------------------
%
\begin{frame}[fragile]{Semantic values}

  \begin{itemize}

    \item The \textbf{semantic values} are defined recursively by the
      following cases:\\
      \smallskip
      \begin{tabular}{@{\,}r@{\;\;}ll}
        $\bullet$
        & \emph{constant}
        & \unit, \textsf{0}, \textsf{1}, \textsf{2}, etc.\\
        $\bullet$
        & \emph{closure}
        & $\clos{x}{e}{\rho}$, with environment $\rho$.\\
        $\bullet$
        & \emph{$n$-tuple}
        & $v_1, \dots, v_n$, with $v_i$ values.
      \end{tabular}

    \item The closure environment $\rho$ is the one at the location
      where the function $\Xfun \; x \rightarrow e$ occurs, and the
      function is closed in that environment (this is called
      \textbf{static scoping}, or \textbf{lexical scoping}).

    \item The evaluation of an expression must take place in the
      context of an environment, where all free variables must be
      bound (in other words, the expression must be closed in that
      environment).

    \item For example, let $\rho$ be the environment before the
      sentence
\begin{alltt}
\small
\textbf{let} x = 1 \textbf{in} x
\end{alltt}
      Then the evaluation of the \emph{second} `x' must done in the
      environment \(\texttt{x} \mapsto 1 \oplus \rho\), where `x' is
      the \emph{first} one.

  \end{itemize}

\end{frame}

% ------------------------------------------------------------------------
%
\begin{frame}{Evaluation}

  Let us proceed to define a semantics for our mini-ML, by associating
  to each expression~\(e\) its value~\(v\), assuming the
  environment~\(\rho\):

  \bigskip

  \begin{tabular}{r@{\,\,}ll@{}}
    $\bullet$
  & $x$
  & Look up the \textbf{last} value bound to~\(x\) in~\(\rho\).\\
    $\bullet$
  & $\Xfun \; x \rightarrow e$
  & The value is the closure $\clos{x}{e}{\rho}$.\\
    $\bullet$
  & $e_1$ \texttt{+} $e_2$, etc.
  & Evaluate $e_1$ and $e_2$ in $\rho$, then add, etc.\\
    $\bullet$
  & \unit \ or \textsf{0} or \textsf{1} or \textsf{2}, etc.
  & The value is \unit or \textsf{0} or \textsf{1}, etc.\\
    $\bullet$
  & $\lpar{e}\rpar$
  & Evaluate $e$ into $v$ in $\rho$.\\
    $\bullet$
  & $\Xlet \; x = e_1 \Xin \; e_2$
  & Evaluate $e_1$ into $v_1$ in $\rho$;\\
  &
  & evaluate $e_2$ into $v$ in $x \mapsto v_1 \oplus \rho$.\\
    $\bullet$
  & $e_1 \; e_2$
  & Evaluate $e_1$ and $e_2$ into $v_1$ and $v_2$ in $\rho$\\
  &
  & (order left unspecified);\\
  &
  & $v_1$ must be the closure $\clos{x}{e}{\rho_1}$;\\
  &
  & evaluate $e$ in $x \mapsto v_2 \oplus \rho_1$:\\
  && the value is $v$.
\end{tabular}

\end{frame}

% ------------------------------------------------------------------------
%
\begin{frame}[fragile]{Example of evaluation}

  \begin{enumerate}

  \item The environment is initially empty. Let us evaluate
\begin{alltt}
\small
\textbf{let} x = 0 \textbf{in}
\textbf{let} id = \textbf{fun} x \(\rightarrow\) x \textbf{in}
\textbf{let} y = id x \textbf{in}
\textbf{let} x = (\textbf{fun} x \(\rightarrow\) \textbf{fun} y \(\rightarrow\) x + y) 1 2
\textbf{in} x + 1
\end{alltt}

  \item The evaluation of
\begin{alltt}
\small
\textbf{let} x = 0 \textbf{in} \ldots
\end{alltt}
    starts with the evaluation of~\texttt{\small 0}, which, obviously,
    yields~\(0\).

   \item We create the binding $\ident{x} \; \mapsto \; 0$, with which
     we extend~\(\rho_\varnothing\), that is, $\ident{x} \mapsto 0
     \oplus \rho_\varnothing$.

   \item We evaluate the elided subexpression within this new
     environment.

  \end{enumerate}

\end{frame}

% ------------------------------------------------------------------------
%
\begin{frame}{Example of evaluation}

  \begin{itemize}

    \item The evaluation of
\begin{alltt}
\small
\textbf{let} id = \textbf{fun} x \(\rightarrow\) x \textbf{in} \ldots
\end{alltt}
      is performed within the environment $\ident{x} \mapsto 0 \oplus
      \rho_\varnothing$.

    \item The value of~\texttt{id} is then the closure
      $\clos{\ident{x}}{\ident{x}}{\ident{x} \mapsto 0 \oplus \rho_\varnothing)}$.

    \item We extend the current environment and evaluate the elided
      subexpression within it.

  \item The evaluation of
\begin{alltt}
\small
\textbf{let} y = id x \textbf{in} \ldots
\end{alltt}
    is done within the environment $\rho = (\ident{id} \mapsto ((\Xfun
    \; \ident{x} \rightarrow \ident{x}), \ident{x} \mapsto 0 \oplus
    \rho_\varnothing)) \oplus (\ident{x} \mapsto 0) \oplus
    \rho_\varnothing$.

    \item We evaluate first (\texttt{\small id x}) in~\(\rho\).

    \item In order to do so, we evaluate~\texttt{\small id}
      and~\texttt{\small x}
      separately.

  \end{itemize}

\end{frame}

% ------------------------------------------------------------------------
%
\begin{frame}{Example of evaluation}

  \begin{itemize}

    \item These are both variables, thus we look them up in the
      environment to retrieve their last corresponding binding:
      \(\rho(\ident{id}) = \clos{\ident{x}}{\ident{x}}{\ident{x}
        \mapsto 0 \oplus \rho_\varnothing}\), and \(\rho(\ident{x}) =
      0\).

    \item We need to evaluate the body in the environment $\rho =
      (\ident{x} \mapsto 0) \oplus (\ident{x} \mapsto 0) \oplus
      \rho_\varnothing$, which yields~\(\rho(\ident{x}) = 0\), etc.

    \item We create the binding $\ident{y} \mapsto 0$, we add it to
      the current environment and we evaluate the elided subexpression
      with the new environment.

    \item Let us stop here and move to expressing the above semantics
      in a formal manner.

  \end{itemize}

\end{frame}

% ------------------------------------------------------------------------
%
\begin{frame}{Formal semantics}

  If we note \(\evalf{e}{\rho}\) the value obtained by evaluating the
  expression~\(e\) in the environment~\(\rho\), then the evaluation of
  mini-ML is:
  \begin{align*}
    \evalf{\overline{n}}{\rho} &= n, \;\, \text{with
      $\overline{n}$ an \OCaml integer and $n \in \mathbb{N}$}\\
    \evalf{e_1 \; \texttt{+} \; e_2}{\rho} &=
    \evalf{e_1}{\rho} + \evalf{e_2}{\rho} \in \mathbb{N}\\
    \evalf{\lpar{e}\rpar}{\rho} &= \evalf{e}{\rho}\\
    \evalf{x}{\rho} &= \rho(x)\\
    \evalf{\Xfun \; x \rightarrow e}{\rho} &= \clos{x}{e}{\rho}\\
    \evalf{e_1 \; e_2}{\rho} &= \evalf{e_3}{(x \mapsto
      \evalf{e_2}{\rho} \oplus \rho_3)},\\
    &\phantom{=\llbracket} \textnormal{where} \ \evalf{e_1}{\rho} =
    \clos{x}{e_3}{\rho_3}\\
    \evalf{\Xlet \; x = e_1 \; \Xin \; e_2}{\rho}
    &= \evalf{(\Xfun \; x \rightarrow e_2) \; e_1}{\rho}.
  \end{align*}

  Note that we write ``$x,\, y$'' instead of ``$(x,\, y)$'' for
  brevity. Evaluation means applying the equations from left to right.

\end{frame}

% ------------------------------------------------------------------------
%
\begin{frame}[fragile]{Partial applications}

  \begin{itemize}

    \item We can easily work out that
      \begin{equation*}
        \evalf{\Xlet \; x = e_1 \; \Xin \; e_2}{\rho}
        = \evalf{e_2}{(x \mapsto \evalf{e_1}{\rho} \oplus \rho)}.
      \end{equation*}

    \item Since closures are values, they can be the value of a
      function call:
\begin{alltt}
\small
\textbf{let} add = \textbf{fun} x \(\rightarrow\) \textbf{fun} y \(\rightarrow\) x + y \textbf{in}
\textbf{let} incr = add 1
\textbf{in} incr 5
\end{alltt}

    \item The call (\texttt{\small add 1}) is a \textbf{partial
      application}, as opposed to a \textbf{complete application} like
      (\texttt{\small add 1 5}). A partial application, by definition,
      always evaluates in a closure.

    \item Operations are functions denoted by a symbol used in infix
      position (between the operands). As such, they can be partially
      evaluated:
\begin{alltt}
\small
\textbf{let} incr = (+) 1 \textbf{in} incr 5
\end{alltt}
     \smallskip
     \noindent where the parentheses in~\texttt{(+)} mean that the
     operator has to be considered in \textbf{prefix} position.

  \end{itemize}

\end{frame}

% ------------------------------------------------------------------------
%
\begin{frame}[fragile]{Auto-application and termination}

  \begin{itemize}

    \item Our mini-ML is \textbf{Turing\hyp{}complete}, just as \OCaml
      is.

    \item For instance, we can write the following
      non\hyp{}terminating program:
      \smallskip
\begin{alltt}
\small
\textbf{let} omega = \textbf{fun} f \(\rightarrow\) f f \textbf{in} omega omega
\end{alltt}

    \item The function \texttt{\small omega} is an
      \textbf{auto\hyp{}applicative function}.

    \item It enables the following non\hyp{}terminating evaluation:
      \begin{gather*}
        \evalf{\texttt{\textbf{let} omega = \textbf{fun} f \(\rightarrow\) f f \textbf{in} omega omega}}{\rho}\\
\begin{align*}
&= \evalf{\texttt{\small omega omega}}{(\texttt{\small omega} \mapsto
    \evalf{\texttt{\small \textbf{fun} f \(\rightarrow\) f f}}{\rho} \oplus
    \rho)}\\
&= \evalf{\texttt{\small f f}}{(\texttt{\small f} \mapsto \clos{f}{f \; f}{\rho})
    \oplus \rho)}\\ &= \text{\emph{idem.}}
\end{align*}
\end{gather*}

  \end{itemize}

\end{frame}

% ------------------------------------------------------------------------
%
\begin{frame}[fragile]{Fixed points}

  \begin{itemize}

    \item To demonstrate the expressive power of our mini-ML, let us
      see how we can simulate recursive functions by means of the
      $\Omega$ function:
\begin{alltt}
\small
\textbf{let} \(\Omega\) = \textbf{fun} f \(\rightarrow\) f f
\end{alltt}

    \item Let us define a function \texttt{\small fix}, called the
      \textbf{Y combinator} in \(\lambda\)-calculus:
      \smallskip
\begin{alltt}
\small
\textbf{let} fix = \textbf{fun} g \(\rightarrow\) \(\Omega\) (\textbf{fun} h \(\rightarrow\) \textbf{fun} x \(\rightarrow\) g (h h) x)
\end{alltt}

      \item For any (terminating) function~\(f\) and
        any function~\(x\),
        \begin{equation*}
          \evalf{f \; \lpar\ident{fix} \; f\rpar{} \; x}{\rho} =
          \evalf{\ident{fix} \; f \; x}{\rho}.
        \end{equation*}

      \item In other words, we have
        \begin{equation*}
          f \; (\ident{fix} \; f) \; x \equiv \ident{fix} \; f \; x.
        \end{equation*}

      \item The fixed point~\(x\) of a (continuous) function~\(g\)
        satisfies \(g(x) = x\).

      \item Therefore, the fixed point of~\ident{f} (if it exists) is
        the value of \lpar\ident{fix} \ident{f}\rpar{}.

  \end{itemize}

\end{frame}

% ------------------------------------------------------------------------
%
\begin{frame}{Fixed points}

   Just as with \Xlet \Xin, we can simplify the semantics of
   \texttt{\small fix}:
   \begin{equation*}
     \evalf{\ident{fix} \; e}{\rho} = \evalf{e_1}{(f \mapsto
       \evalf{\ident{fix} \, (\Xfun \, f \rightarrow e_1)}{\rho}
       \oplus \rho_1)},
   \end{equation*}
   where $\evalf{e}{\rho} = \clos{f}{e_1}{\rho_1}$.

   \bigskip

   There is a choice:
   \begin{itemize}

     \item either we define \texttt{\small fix} in the functional
       language (as a built\hyp{}in constant), or

     \item at the meta\hyp{}language (the mathematical logic we are
       using to define the semantics).

   \end{itemize}

\end{frame}

% ------------------------------------------------------------------------
%
\begin{frame}[fragile]{Fixed points}

  Let us consider the following definitions:
\begin{alltt}
\small
\textbf{let} pre\_fact =
  \textbf{fun} f \(\rightarrow\) \textbf{fun} n \(\rightarrow\) \textbf{if} n = 1 \textbf{then} 1 \textbf{else} n * f(n-1)
\textbf{let} fact = fix pre\_fact
\end{alltt}

  We see that \texttt{\small fact} is the fixed point of
  \texttt{\small pre\_fact}:
  \begin{align*}
    \evalf{\texttt{\small fact}}{\rho} &= \evalf{\texttt{\small fix
        pre\_fact}}{\rho}\\
    &= \evalf{\texttt{\small pre\_fact (fix pre\_fact)}}{\rho}\\
    &= \evalf{\texttt{\small pre\_fact fact}}{\rho}\\
    &= \evalf{\texttt{\small \textbf{fun} n \(\rightarrow\) \textbf{if} n
        = 1 \textbf{then} 1 \textbf{else} n * fact(n-1)}}{\rho}
  \end{align*}

  It is as if the definition of `\texttt{\small fact}' were
  recursive... but it is not.

\end{frame}

% ------------------------------------------------------------------------
%
\begin{frame}[fragile]{Fixed points in C}

\begin{alltt}
\small
\textbf{#include}<stdio.h>
\textbf{#include}<stdlib.h>

\textbf{typedef int} (*fp)();

\textbf{int} fact(fp f, \textbf{int} n) \{
  \textbf{return} n? n * ((\textbf{int} (*)(fp,\textbf{int}))f)(f,n-1) : 1; \}

\textbf{int} read(\textbf{int} dec, \textbf{char} arg[]) \{
  \textbf{return} ('0' <= *arg && *arg <= '9')?
         read(10*dec + *arg - '0', arg+1) : dec; \}

\textbf{int main}(\textbf{int} argc, \textbf{char}** argv) \{
  \textbf{if} (argc == 2)
     printf("%u\textbackslash{n}", fact(&fact, read(0, argv[1])));
  \textbf{else} printf("Only one integer allowed.\textbackslash{n}");
  \textbf{return} 0; \}
\end{alltt}
\end{frame}

%% % ------------------------------------------------------------------------
%% %
%% \begin{frame}{Local recursive definitions}

%%   \begin{itemize}

%%     \item The function `\texttt{\small fact}' coincides pointwise with
%%       the factorial function.

%%     \item Let us extend mini-ML with a native recursive binder, using
%%       the `\texttt{\small fix}' function:
%%       \begin{equation*}
%%         \evalf{\Xlet \; \Xrec \; f = e_1 \; \Xin \; e_2}{\rho} =
%%         \evalf{\Xlet \; f = \textsf{fix} \; (\Xfun \; f \rightarrow
%%           e_1) \; \Xin \; e_2}{\rho}.
%%       \end{equation*}

%%   \end{itemize}

%% \end{frame}

%% % ------------------------------------------------------------------------
%% %
%% \begin{frame}{Extensions}

%%   \begin{itemize}

%%     \item Let us extend our mini-ML with the following
%%       expressions.\\
%%       \smallskip
%%       \begin{tabular}{r@{\,\,}ll@{}}
%%         $\bullet$
%%         & \emph{Boolean constants}
%%         & \Xtrue{} or \Xfalse\\
%%         $\bullet$
%%         & \emph{Boolean operators}
%%         & \texttt{(\&\&)} or \texttt{(||)} or \textsf{\textbf{not}}\\
%%         $\bullet$
%%         & \emph{$n$-tuple}
%%         & $e_1, \ldots, e_n$\\
%%         $\bullet$
%%         & \emph{conditional}
%%         & \Xif{} $e_0$ \Xthen{} $e_1$ \Xelse{} $e_2$\\
%%         $\bullet$
%%         & \emph{local recursive binding}
%%         & \Xlet{} \Xrec{} $f \;\; = \;\; e_1$ \Xin{} $e_2$
%%       \end{tabular}

%%       \item Note that parentheses are recommended around tuples, but
%%         are not mandatory.

%%       \item Let us distinguish the variables occurring after \Xlet{}
%%         and \Xfun{}, because they are \textbf{irrefutable patterns},
%%         noted $\ipat{p}$:\\
%%         \smallskip
%%         \begin{tabular}{r@{\,\,}ll@{}}
%%           $\bullet$
%%           & \emph{function}
%%           & $\Xfun \;\; \ipat{p} \rightarrow e$\\
%%           $\bullet$
%%           & \emph{local definition}
%%           & \(\Xlet \;\; \textrm{[}\Xrec\textrm{]} \;\; \ipat{p} \;\;
%%           = \;\; e_1 \;\; \Xin \;\; e_2\)
%%         \end{tabular}

%%   \end{itemize}

%% \end{frame}

%% % ------------------------------------------------------------------------
%% %
%% \begin{frame}{Irrefutable patterns}

%%   \begin{itemize}

%%     \item An irrefutable pattern $\ipat{p}$ is defined recursively by
%%       the following cases:\\
%%       \smallskip
%%       \begin{tabular}{r@{\,\,}ll@{}}
%%         $\bullet$
%%         & \emph{variable}
%%         & $f, g, h, x, y, z$, etc. \\
%%         $\bullet$
%%         & \emph{unit}
%%         & \unit\\
%%         $\bullet$
%%         & \emph{$n$-tuple}
%%         & $\ipat{p}_1, \ldots, \ipat{p}_n$\\
%%         $\bullet$
%%         & \emph{parenthesis}
%%         & $\lpar\ipat{p}\rpar$\\
%%         $\bullet$
%%         & \emph{underscore}
%%         & {\Large \_}
%%       \end{tabular}

%%       \item Note that, from the syntactic point of view, irrefutable
%%         patterns are special cases of expressions, except the
%%         underscore, which is a special case avoiding the creation of a
%%         binding.

%%       \item The sentence~\(e\) is equivalent to \(\Xlet \;\; \_ \;\; =
%%         \;\; e\).

%%   \end{itemize}

%% \end{frame}
